% TOP_PROBLEM: {"id": 8700033, "title": "Conjecture 3.15 — Goncharov", "category_id": 3, "category": {"id": 3, "name": "graph_theory", "display_name": "Graph Theory", "description": "Problems involving graphs, networks, and their properties.", "slug": "graph-theory", "order_index": 3, "created_at": "2026-07-31T15:26:25.671Z"}, "status": "partially_solved", "problem_number": "AMR-086-0033", "set_id": 13}
% FIRST_POSTED: 2026-10-01
% ENTRY_KIND: statement_only
% UnsolvedMath ID 8700033; source catalogue code AMR-086-0033
% !TeX program = lualatex
% Definition and sources only; this file is not an LLM solution attempt.
\documentclass[11pt,a4paper]{article}
\usepackage[margin=25mm]{geometry}
\usepackage{fontspec}
\setmainfont{lmroman10-regular.otf}[BoldFont=lmroman10-bold.otf,
 ItalicFont=lmroman10-italic.otf,BoldItalicFont=lmroman10-bolditalic.otf]
\usepackage{amsmath,amssymb,mathtools}
\usepackage{unicode-math}
\setmathfont{latinmodern-math.otf}
\AtBeginDocument{\let\setminus\smallsetminus}
\usepackage{xurl,hyperref}
\hypersetup{colorlinks=true,urlcolor=blue}
\setcounter{secnumdepth}{0}
\setlength{\parindent}{0pt}
\setlength{\parskip}{5pt}
\setlength{\emergencystretch}{3em}
\begin{document}
\section*{Conjecture 3.15 — Goncharov}\label{8700033}
{\small UnsolvedMath ID 8700033 · AMR-086-0033 · Graph Theory · AMR Open Problem Lists}
\subsection{Definitions and mathematical statement}
For integers $r\ge1$, $s_1\ge2$, and $s_2,\ldots,s_r\ge1$, define the convergent multiple zeta value
\[
\zeta(s_1,\ldots,s_r)=
\sum_{n_1>\cdots>n_r\ge1}\frac1{n_1^{s_1}\cdots n_r^{s_r}}.
\]
Its weight is $s_1+\cdots+s_r$. For $w\ge2$, let $\mathfrak Z_w$ be the rational vector space spanned by all these values of weight $w$. Set $\mathfrak Z_0=\mathbb Q$, $\mathfrak Z_1=\{0\}$, and $\mathfrak Z=\sum_{w\ge0}\mathfrak Z_w\subset\mathbb R$, where every sum is finite.

Products of multiple zeta values expand as rational linear combinations of values with the combined weight, so $\mathfrak Z_u\mathfrak Z_v\subseteq\mathfrak Z_{u+v}$. Is the sum defining $\mathfrak Z$ direct? Precisely, must every finite relation
\[
z_0+z_1+\cdots+z_W=0,\qquad z_w\in\mathfrak Z_w,
\]
have $z_w=0$ for each weight $w$?

This would make the actual algebra of these real numbers graded by weight. The assertion is about numerical values, rather than formal symbols assigned a weight by definition. Relations within one weight are allowed, and no dimension formula for an individual space is asserted here.

Multiplication converts polynomial relations into finite sums organized by total weight. Thus the conjecture says that cancellation between genuinely different weights cannot create an additional relation. The admissibility convention $s_1\ge2$ ensures convergence and avoids requiring regularized divergent zeta values in the definition.
\subsection{Short English statement}
Do numerical multiple zeta values admit a genuine grading by weight, with no nonzero relation mixing different weights?
\subsection{Sources}
\begin{itemize}
\item[S1] ULAM.AI and the UnsolvedMath Contributors, supplied problems.json, record 8700033 (AMR-086-0033), AMR Open Problem Lists. Catalogue identity and source transcription; CC BY 4.0.
\url{https://huggingface.co/datasets/ulamai/UnsolvedMath}.
\item[S2] Waldschmidt - Open Diophantine Problems (2004), Conjecture 3.15. Original source indexed by AMR-086-0033.
\url{https://arxiv.org/abs/math/0312440}.
\item[S3] M. Waldschmidt, Open Diophantine Problems, Moscow Mathematical Journal 4 (2004), 245–305. Author PDF supplies the surrounding definitions and hypotheses.
\url{https://webusers.imj-prg.fr/~michel.waldschmidt/articles/pdf/odp.pdf}.
\end{itemize}
\end{document}
